\section{Análisis de consenso\_prom.launch.py}

\subsection{Visión general}

consenso\_prom.launch.py es el launch file principal del experimento de consenso de promedio. Orquesta el arranque de Gazebo, $n$ robots (hasta 8), el nodo de conectividad AIRE, y dos nodos de graficación. Todo cabe en un único archivo gracias a la composición de los patrones descritos en las secciones anteriores.

La estructura del archivo sigue el patrón OpaqueFunction: generate\_launch\_description declara los argumentos y delega toda la lógica de construcción a launch\_setup. Esto es necesario porque el número de nodos a lanzar depende del argumento n\_robots, que no se conoce hasta tiempo de ejecución.

\subsection{Argumentos declarados}

El Código~\ref{lst:consenso_args} muestra los seis argumentos que el usuario puede sobreescribir desde la terminal:

\begin{figure*}[t]
\centering
\begin{lstlisting}[style=python,
    caption={Argumentos de consenso\_prom.launch.py.},
    label={lst:consenso_args}]
DeclareLaunchArgument('n_robots',         default_value='8'),
DeclareLaunchArgument('neighbor_radius',  default_value='5.0'),
DeclareLaunchArgument('t_final',          default_value='60.0'),
DeclareLaunchArgument('spawn_obstacle',   default_value='false'),
DeclareLaunchArgument('obstacle_x',       default_value='4.0'),
DeclareLaunchArgument('obstacle_y',       default_value='4.5'),
\end{lstlisting}
\end{figure*}


Un experimento típico con 4 robots y 30 segundos se lanza con el comando de la Terminal~\ref{lst:consenso_launch_cmd}:

{\renewcommand{\lstlistingname}{Terminal}
\begin{figure*}[t]
\centering
\begin{lstlisting}[style=bash,
    caption={Lanzamiento del experimento con argumentos personalizados.},
    label={lst:consenso_launch_cmd}]
ros2 launch simulador consenso_prom.launch.py n_robots:=4 t_final:=30.0
\end{lstlisting}
\end{figure*}
}


\subsection{Poses iniciales predefinidas}

Las posiciones y orientaciones iniciales de los 8 robots están hardcodeadas como una constante POSES\_8 al inicio del archivo. El diseño del layout garantiza conectividad con radio R = 5.0 m: ningún par de robots está a más de 5 m de distancia de todos sus vecinos. Al usar solo n\_robots de los 8 disponibles se toman los primeros n de la lista, lo que permite ejecutar experimentos con subconjuntos del layout completo sin rediseñar las poses.

\subsection{Secuencia de lanzamiento}

La función launch\_setup construye la lista actions en el siguiente orden:

\textbf{1. Gazebo.} La primera acción (sin timer) es la inclusión de gazebo.launch.py. Arranca inmediatamente, dando el mayor tiempo posible para que el simulador se inicialice antes de que lleguen los primeros spawns.

\textbf{2. Obstáculo opcional.} Si spawn\_obstacle = true, se agrega un TimerAction a 3.0 s que spawnea un cilindro en la posición indicada. El retardo de 3 s da tiempo a Gazebo para estar listo.

\textbf{3. Robots (bucle).} Para cada robot $i$, con retardo base delay = 3.0 + i*1.5:
\begin{itemize}
  \item RSP en delay s.
  \item Spawn en delay + 0.5 s.
  \item Nodo de control en delay + 1.5 s.
\end{itemize}

\textbf{4. Nodos tardíos.} Una vez que el último robot está operativo (late\_delay = 3.0 + n*1.5 + 2.0), se lanzan AIRE y los dos plotters.

El Diagrama~\ref{lst:timeline} ilustra la línea de tiempo para n\_robots=3:

{\renewcommand{\lstlistingname}{Diagrama}
\begin{figure*}[t]
\centering
\begin{lstlisting}[style=bash,
    caption={Línea de tiempo de arranque para n\_robots=3.},
    label={lst:timeline}]
t=0.0 s  -> Gazebo arranca
t=3.0 s  -> RSP robot0 | spawn_obstacle (si aplica)
t=3.5 s  -> spawn robot0
t=4.5 s  -> control robot0 | RSP robot1
t=5.0 s  -> spawn robot1
t=6.0 s  -> control robot1 | RSP robot2
t=6.5 s  -> spawn robot2
t=7.5 s  -> control robot2
t=9.5 s  -> AIRE + plotters  (3.0 + 3*1.5 + 2.0)
\end{lstlisting}
\end{figure*}
}


\subsection{Parámetros del nodo de control}

El nodo consenso\_node recibe sus parámetros directamente en el launch file, sin archivo YAML. Esto es práctico para experimentos: los valores pueden cambiarse sin recompilar ni mantener archivos de configuración adicionales. El Código~\ref{lst:consenso_params} muestra los parámetros configurados:

\begin{figure*}[t]
\centering
\begin{lstlisting}[style=python,
    caption={Parámetros del nodo de consenso en el launch file.},
    label={lst:consenso_params}]
control_node = Node(
    package='control_nodes',
    executable='consenso_node',
    name='consenso',
    namespace=ns,
    parameters=[{
        'consensus_type': 'prom_err',
        'k1':   0.8,
        'k2':   1.0,
        'vmax': 1.0,
        'wmax': 1.0,
    }],
    output='screen'
)
\end{lstlisting}
\end{figure*}

